\section{An integer pencil and its exact certificate}
\label{sec:certificate}

This section verifies every hypothesis of Proposition~\ref{prop:tensor-criterion}
for an explicit map $L:\Mat4\to\Mat4$.  Four integer $6$-by-$4$ matrices
$M_0,\ldots,M_3$ define $L$ and the pencil $M(x)=\sum_{i=0}^3x_iM_i$.
The maximal minors of this pencil identify all singular projective directions.
Exact reductions at three primes, together with the quotient and Galois
arguments below, control quadratic evaluations on every ten of those
directions.

\subsection{The explicit matrix pencil}

Use indices $0,1,2,3$ and define real $6$-by-$4$ matrices
\begin{equation}
\begin{aligned}
M_0&=\begin{pmatrix}
6&0&0&0\\0&6&0&0\\0&0&6&0\\0&0&0&6\\0&0&0&0\\0&0&0&0
 \end{pmatrix},&
M_1&=\begin{pmatrix}
-12&0&0&0\\0&-6&0&0\\0&0&6&0\\0&0&0&12\\-6&0&6&6\\-6&-6&6&0
 \end{pmatrix},\\[2mm]
M_2&=\begin{pmatrix}
0&6&-2&0\\6&6&0&2\\-2&0&10&0\\0&2&0&24\\0&0&0&6\\0&-6&6&-6
 \end{pmatrix},&
M_3&=\begin{pmatrix}
0&0&-4&-3\\0&-2&-3&-4\\-4&-3&11&6\\-3&-4&6&25\\6&6&0&0\\0&-6&6&6
 \end{pmatrix}.
\end{aligned}
\label{eq:explicit-matrices}
\end{equation}
For $x=(x_0,x_1,x_2,x_3)\in\C^4$, put $M(x)=\sum_{i=0}^3x_iM_i$ and define
the complex-linear map $L$ on matrix units by
\begin{equation}
 L(E_{ij})=M_i^{\T}M_j\qquad(0\leq i,j<4).
 \label{eq:L}
\end{equation}

\begin{lemma}\label{lem:pencil-ppt}
The map $L:\Mat4\to\Mat4$ in \eqref{eq:L} is PPT.  For every $x\in\C^4$,
\begin{equation}
 L(P_x)=M(x)^{\T}\overline{M(x)}=M(x)^*M(x),\qquad
 \rank L(P_x)=\rank M(x).
 \label{eq:gram-identity}
\end{equation}
Moreover, $L(P_{e_0})=36I_4$.
\end{lemma}

\begin{proof}
Direct integer multiplication gives
\begin{equation}
 M_i^{\T}M_j=M_j^{\T}M_i\qquad(0\leq i,j<4).
 \label{eq:block-symmetry}
\end{equation}
The exact checker in Appendix~\ref{app:arithmetic} verifies these
equalities.  For each row $r$, let $Q_r$ be the $4$-by-$4$ matrix whose
$i$th column is the transpose of row $r$ of $M_i$.  Comparing entries gives
$L(A)=\sum_{r=1}^6Q_rAQ_r^*$, so $L$ is CP.  Equation
\eqref{eq:block-symmetry} also makes every $M_i^{\T}M_j$ symmetric.
Therefore $T_4\circ L=L$, proving the PPT property.

For complex $x$, expansion of \eqref{eq:L} gives
\[
 L(P_x)=\sum_{i,j}x_i\overline{x_j}M_i^{\T}M_j
       =M(x)^{\T}\overline{M(x)}.
\]
Interchanging $i,j$ and using \eqref{eq:block-symmetry} gives the second
expression in \eqref{eq:gram-identity}.  Since
$v^*M(x)^*M(x)v=\norm{M(x)v}^2$, the kernel of this Gram matrix is
$\ker M(x)$, which proves the rank equality.  Finally,
$M_0^{\T}M_0=36I_4$.
\end{proof}

The positive-definite input required by the tensor criterion is now
available.  The remaining assertion concerns every complex rank-loss
direction of the pencil.

\begin{proposition}[Twenty rank-loss directions with independent quadratic evaluations]
\label{prop:twenty-directions}
The projective directions $[x]\in\mathbb P^3(\C)$ for which
$\rank M(x)<4$ consist of exactly twenty distinct points.  Each has a
unique representative $x=(1,t_1,t_2,t_3)$.  Every nonzero homogeneous
quadratic in $\C[x_0,x_1,x_2,x_3]$ vanishes on at most nine of these
representatives.
\end{proposition}

The equation $M(x)y=0$ with $x,y\ne0$ is related to the bilinear kernel
geometry studied by Hansen, Hauge, Myrheim, and Sollid for $4$-by-$4$ PPT
states whose matrix and partial transpose both have rank six
\cite[Sections 2.2 and 7]{HansenEtAl2012}.
Their generic count concerns twenty product rays $x\otimes y$.  The
quadratic condition here concerns the vectors $x\otimes x$ for this
specific pencil, and the following argument proves it from exact data.

\noindent\textit{Proof of Proposition~\ref{prop:twenty-directions}.}
The maximal-minor relations first show that the affine quotient is nonzero
and spanned by twenty monomials.  We then construct a formal multiplication
matrix.  Its irreducible characteristic polynomial, certified modulo $41$,
proves that the spanning monomials are independent and that the quotient
is a degree-twenty field.  This identifies all twenty complex directions.
Finally, reductions at $131$ and $139$ give the full symmetric Galois action,
which transports one independent set of ten quadratic evaluations to every
such set.  The three general algebraic facts used along the way are stated
at their uses; their proofs are collected in
Appendix~\ref{app:certificate-algebra}.

The checker in Appendix~\ref{app:arithmetic} supplies the integer block
identities and the finite modular data used in these stages.  It does not
enumerate the twenty complex directions or their ten-point determinants.
The quotient, projective, and Galois deductions are proved below.

\subsection{The maximal minors and a nonzero quotient}

Because $M(x)$ has four columns, the condition $\rank M(x)<4$ is equivalent
to simultaneous vanishing of its fifteen maximal minors.  For each
four-element subset $\rho\subset\{1,\ldots,6\}$, taken in lexicographic
order with its rows in increasing order, let $\Delta_\rho(x)$ be the
determinant of those rows of $M(x)$.  Put
$p_\rho(t)=\Delta_\rho(1,t_1,t_2,t_3)$ and define
\[
 I=(p_\rho)_\rho\subset\Q[t_1,t_2,t_3],\qquad
 \mathcal A=\Q[t_1,t_2,t_3]/I.
\]
These equations describe the singular directions in the chart $x_0=1$.
We have not yet proved that $I$ is proper, or that this chart contains all
the singular directions.

Order monomials first by increasing total degree $d$ and, within degree
$d$, by the exponent triples
\begin{equation}
 (a,b,d-a-b),\qquad a=0,\ldots,d,\quad b=0,\ldots,d-a,
 \label{eq:monomial-order}
\end{equation}
in the displayed nested order.  Let $v$ be the column of the twenty
monomials of degree at most three, and let $w$ be the column of the fifteen
monomials of degree four.  The coefficients of the minors define integer
matrices $D,H$ by
\begin{equation}
 (p_\rho(t))_\rho=Dv(t)+Hw(t),\qquad
 D\in\mathbb Z^{15\times20},\quad H\in\mathbb Z^{15\times15}.
 \label{eq:DH}
\end{equation}

The exact coefficient calculation in Appendix~\ref{app:arithmetic} gives
$\det H\equiv40\pmod{41}$, so $H$ is invertible over $\Q$.  Define
$C=-H^{-1}D$.  In the quotient $\mathcal A$, Equation~\eqref{eq:DH} gives
the relations $w=Cv$, expressing each degree-four monomial as a linear
combination of the twenty lower monomials.

The relations $w=Cv$ must describe a nonzero quotient before they can
be used to count its points.  We first verify this and show that the low
monomials span.  The existence argument uses the following intersection
fact for $\mathbb P^3\times\mathbb P^3$.

\begin{lemma}[Six bilinear equations]\label{lem:six-bilinear}
Any six complex bilinear forms on $\C^4\times\C^4$ have a common zero
$(x,y)$ with $x\ne0$ and $y\ne0$.
\end{lemma}

The proof is given in Appendix~\ref{app:certificate-algebra}.

Apply the lemma to the six rows of
$M(x)y$.  It gives $x,y\ne0$ with $M(x)y=0$, so all $\Delta_\rho(x)$
vanish.

The homogeneous degree-four part of $p_\rho(t)$ is
$\Delta_\rho(0,t_1,t_2,t_3)$.  Equation~\eqref{eq:DH} therefore gives
\[
 \bigl(\Delta_\rho(0,t)\bigr)_\rho=Hw(t).
\]
If a common zero had $x_0=0$, invertibility of $H$ would imply $w(t)=0$.
The entries $t_1^4,t_2^4,t_3^4$ occur in $w$, so $t=0$, contradicting
$x\ne0$.  Thus every projective common zero lies in the chart $x_0=1$.
Rescaling the zero furnished by the lemma gives an affine common zero of
the $p_\rho$, proving that $I$ is proper and $\mathcal A\ne0$.  The same
argument has excluded all points at infinity.

In $\mathcal A$, the relations $w=Cv$ show that the twenty low monomials
span.  In a monomial of total degree at least four, substitute for any
degree-four factor.  Each resulting term has strictly smaller total
degree.  Induction reduces every monomial to a linear combination of
the low monomials; this proves spanning without assuming unique
reductions.

\subsection{A formal multiplication operator}

Multiplication matrices are a standard way to encode finite polynomial
systems; see Sottile \cite[Section 2.3]{Sottile2001}.  The quotient
$\mathcal A$ is now known to be nonzero and spanned by the
twenty low monomials, but their independence remains to be proved.
Take a formal twenty-dimensional rational vector space with basis
$e_\alpha$ indexed by exponent triples $|\alpha|\leq3$, in the order
\eqref{eq:monomial-order}.  Define an endomorphism $N$ of this formal
space and its characteristic polynomial by
\begin{equation}
 Ne_\alpha=
 \begin{cases}
 e_{\alpha+(1,0,0)},&|\alpha|\leq2,\\[1mm]
 \displaystyle\sum_{|\beta|\leq3}
 C_{\alpha+(1,0,0),\beta}e_\beta,&|\alpha|=3,
 \end{cases}
 \qquad f(z)=\det(zI_{20}-N).
 \label{eq:N}
\end{equation}
The row index of $C$ in the second line is the indicated degree-four
monomial.  Let $\pi:\Q^{20}\to\mathcal A$ send $e_\alpha$ to the class of
$t^\alpha$, and let $\mu_1$ be multiplication by the class of $t_1$.
The spanning result makes $\pi$ surjective.  By the definition of $N$,
\begin{equation}
 \pi N=\mu_1\pi.
 \label{eq:intertwining}
\end{equation}
In particular, $\ker\pi$ is an $N$-invariant rational subspace.
We will use irreducibility of $f$ to show that this kernel is zero,
which will identify $N$ with multiplication by $t_1$.

For a prime $p$ with $\det H\not\equiv0\pmod p$, set
\[
 \mathbb Z_{(p)}=\{a/b:a,b\in\mathbb Z,\ p\nmid b\}.
\]
The identity $H^{-1}=\operatorname{adj}(H)/\det H$ shows that the entries
of $C,N$ and the coefficients of the monic polynomial $f$ lie in
$\mathbb Z_{(p)}$.  In particular their denominators are units at $p$,
and reduction gives
$\overline f=\det(zI_{20}-\overline N)\in\F_p[z]$.  This is a statement
of integrality at the specified prime; no global integrality of the
coefficients of $f$ is needed.

For $k\geq1$, put
\[
 d_p(k)=\deg\gcd(\overline f,z^{p^k}-z).
\]
The exact calculation in Appendix~\ref{app:arithmetic} gives the computed
columns of Table~\ref{tab:modular-certificate}.  Its output also records all twenty-one characteristic
coefficients and all twenty gcd degrees at each prime.

\begin{table}[htbp]
\centering
\caption{Computed modular data and the factor degrees deduced from them.}
\label{tab:modular-certificate}
\begin{tabular}{@{}ccll@{}}
\toprule
$p$ & $\det H\bmod p$ & \shortstack{computed\\gcd degrees} &
\shortstack{deduced\\factor degrees}\\
\midrule
41 & 40 & $d_{41}(k)=0\ (1\leq k<20),\ d_{41}(20)=20$ & $20$\\
131 & 23 & $d_{131}(1)=1,\ d_{131}(2)=3,\ d_{131}(17)=18$ & $1,2,17$\\
139 & 9 & $d_{139}(1)=1,\ d_{139}(19)=20$ & $1,19$\\
\bottomrule
\end{tabular}
\end{table}

The determinant residues justify all three local reductions.  The primes
exceed twenty, so the divisions in the Newton identities used to compute
the characteristic polynomials are valid.  Their primality follows by
trial division by the primes at most eleven.

We explain the last column of the table.  Over $\F_p$ the polynomial
$z^{p^k}-z$ is squarefree, since its derivative is $-1$.  Its irreducible
factors are precisely the monic irreducible polynomials whose degrees
divide $k$: a root of an irreducible polynomial of degree $s$ has a
Frobenius orbit of length $s$.  Hence $d_p(k)$ is the sum of the degrees of
the \emph{distinct} irreducible factors of $\overline f$ whose degrees
divide $k$.  At $p=41$, a factor of degree $s<20$ would make $d_{41}(s)$
positive, so $\overline f$ is irreducible of degree twenty.  At $p=131$,
the three displayed values force one factor of degree one, one of degree
two, and one of degree seventeen.  These degrees already sum to twenty,
so there are no further or repeated factors.  The two values at $p=139$
similarly force one factor of degree one and one of degree nineteen,
without repetition.  All three reductions are squarefree.

The following elementary reduction lemma lets us infer rational
irreducibility without assuming that $f$ has integer coefficients.
Its $p$-adic form is used in the proof of Lemma~\ref{lem:local-frobenius}
in Appendix~\ref{app:certificate-algebra}.

\begin{lemma}[Monic local reduction]\label{lem:monic-local-reduction}
Let $p$ be a prime, and let the coefficient field and its valuation ring be
one of the pairs
\[
 (\mathbb L,R_p)=(\Q,\mathbb Z_{(p)})
 \quad\text{or}\quad
 (\mathbb L,R_p)=(\Q_p,\mathbb Z_p).
\]
If a monic polynomial $P\in R_p[z]$ factors as $P=gh$ with monic
$g,h\in\mathbb L[z]$, then $g,h\in R_p[z]$, and their reductions modulo
$p$ retain their degrees, including degree zero.  Consequently, if $P$
has positive degree and irreducible reduction in $\F_p[z]$, then $P$ is
irreducible over $\mathbb L$.
\end{lemma}

The proof is given in Appendix~\ref{app:certificate-algebra}.

Applying the lemma to $f\in\mathbb Z_{(41)}[z]$ and its irreducible
reduction proves that $f$ is irreducible over $\Q$.

\subsection{The coordinate field and its twenty points}

We can now prove independence of the twenty spanning monomials.
An endomorphism
with irreducible characteristic polynomial has no nonzero proper invariant
subspace: a basis adapted to such a subspace makes the matrix block
triangular and factors its characteristic polynomial into those of the
restriction and quotient.  Since $\pi$ is surjective and
$\mathcal A\ne0$, its kernel is proper.  Irreducibility of $f$ forces
$\ker\pi=0$.  The twenty low monomials are therefore a basis of
$\mathcal A$, and $N$ is its multiplication matrix for $t_1$.

By Cayley--Hamilton the minimal polynomial of $N$ divides $f$.  It is
nonconstant, so irreducibility of $f$ makes it equal to $f$.  In a nonzero
unital algebra, a polynomial in an element vanishes exactly when the
same polynomial in its multiplication operator vanishes; the reverse
implication follows by applying the operator to $1$.  Hence the
subalgebra generated by $t_1$ is $\Q[z]/(f)$ and has dimension twenty.
It is already all of $\mathcal A$:
\begin{equation}
 \mathcal A\simeq\Q[z]/(f),\qquad z\longmapsto t_1.
 \label{eq:coordinate-field}
\end{equation}
In particular $\mathcal A$ is a field.  Write the classes of $t_2,t_3$
as $h_2(t_1),h_3(t_1)$ with $h_2,h_3\in\Q[z]$ of degree below twenty.

The irreducible polynomial $f$ is separable in characteristic zero.
Let its distinct complex roots be $\lambda_1,\ldots,\lambda_{20}$.
Extending scalars in \eqref{eq:coordinate-field} and applying the Chinese
remainder theorem gives
\begin{equation}
 \C[t_1,t_2,t_3]/I\C[t_1,t_2,t_3]
 \simeq\C[z]/(f)\simeq\prod_{j=1}^{20}\C.
 \label{eq:split-algebra}
\end{equation}
The factors are evaluations at the distinct triples
\begin{equation}
 t^{(j)}=(\lambda_j,h_2(\lambda_j),h_3(\lambda_j)).
 \label{eq:points}
\end{equation}
Conversely, any complex common zero of the $p_\rho$ defines a unital
$\C$-algebra homomorphism from the left side of
\eqref{eq:split-algebra} to $\C$, and such a homomorphism selects one
factor of the product.  Thus the triples in \eqref{eq:points} are all the
affine common zeros.  The product of fields also shows that the complex
ideal is radical, so these twenty points are reduced.  Together with the
exclusion of $x_0=0$, this identifies all complex projective singular
directions of the pencil.


\subsection{Local Frobenius and quadratic evaluations}

It remains to prove that no nonzero quadratic vanishes on ten of the
directions.  We will show that the Galois group of the splitting field
permutes the twenty roots arbitrarily.  This transports independence for
one set of ten evaluations to every set of ten.

The following is the locally integral form of the Frobenius
factorization principle, often called Dedekind's theorem.  Milne states
the monic integer case \cite[Theorem 4.28]{Milne2022}.  We give the local
rational argument needed for the coefficients of $f$.

\begin{lemma}[Locally integral Frobenius]\label{lem:local-frobenius}
Let $F\in\Q[z]$ be monic, separable, and of positive degree, with
coefficients in $\mathbb Z_{(p)}$.  If its reduction in $\F_p[z]$ is
squarefree with irreducible-factor degrees $s_1,\ldots,s_h$, then the
Galois group of its rational splitting field contains an element with
cycle lengths $s_1,\ldots,s_h$ on the roots.
\end{lemma}

The proof is given in Appendix~\ref{app:certificate-algebra}.

Let $\mathbb K\subset\C$ be a splitting field of $f$, and let
$G=\operatorname{Gal}(\mathbb K/\Q)$ act on its twenty roots.  Applying
Lemma~\ref{lem:local-frobenius} to the reductions at $131$ and $139$
gives elements of cycle types $(1,2,17)$ and $(1,19)$.  Irreducibility
of $f$ makes this action transitive.  A nineteen-cycle fixes one root and
is transitive on the other nineteen, so the stabilizer of that root is
transitive on its complement.  Transitivity conjugates this property to
every point stabilizer; hence $G$ is transitive on ordered pairs of
distinct roots, or two-transitive.  The seventeenth power of an element
of type $(1,2,17)$
is a transposition.  Two-transitivity sends its two entries to any
ordered pair, and conjugation therefore gives every transposition.  Thus
\begin{equation}
 G=S_{20}.
 \label{eq:galois-group}
\end{equation}
This is the elementary permutation criterion recorded in
Milne \cite[Lemma 4.31]{Milne2022}.  The group in
\eqref{eq:galois-group} is the splitting-field group acting on the
twenty roots.

Form the $20$-by-$10$ evaluation matrix over $\mathbb K$
\[
 \mathcal E_{j,\alpha}=(t^{(j)})^\alpha,\qquad
 1\leq j\leq20,\quad |\alpha|\leq2.
\]
Its entries lie in $\mathbb K$ by \eqref{eq:points}.  The ten column
monomials belong to the basis of $\mathcal A$ already proved above.
They remain independent after scalar extension to $\mathbb K$, where
the split algebra is $\mathbb K^{20}$ through evaluation at the
twenty triples.  Hence $\mathcal E$ has column rank ten over
$\mathbb K$, so some ten-row square submatrix has nonzero determinant.

Every $\sigma\in G$ permutes the triples in \eqref{eq:points} according
to its permutation of the roots, since $h_2,h_3$ have rational
coefficients.  In particular
$\sigma(\mathcal E_{j,\alpha})=\mathcal E_{\sigma(j),\alpha}$.
The group $S_{20}$ is transitive on ten-element subsets.  Applying a
suitable $\sigma$ to the known nonzero determinant therefore gives,
up to a row-order sign, the determinant for any prescribed ten rows.
A field automorphism preserves nonzero elements, so every such
determinant is nonzero in $\mathbb K$ and hence in $\C$.  The Galois
action has been used only on algebraic determinants; their nonvanishing
gives invertibility for arbitrary complex coefficient vectors.

Every homogeneous quadratic has a unique expression
\[
 q(x_0,x_1,x_2,x_3)=\sum_{|\alpha|\leq2}
 a_\alpha x_0^{2-|\alpha|}
 x_1^{\alpha_1}x_2^{\alpha_2}x_3^{\alpha_3}.
\]
Its dehomogenization $q(1,t)$ is nonzero exactly when $q$ is nonzero.
If it vanished at ten triples, the corresponding invertible evaluation
matrix would force every $a_\alpha=0$.  Thus it vanishes on at most
nine triples.  The identification of all common zeros above, the
exclusion of $x_0=0$, and the equivalence between vanishing of the
maximal minors and $\rank M(x)<4$ now prove
Proposition~\ref{prop:twenty-directions}.\qed

\begin{proof}[Proof of Theorem~\ref{thm:main-pair}]
Take $L$ from \eqref{eq:L}.  Lemma~\ref{lem:pencil-ppt} supplies its PPT
property, the rank identity, and the positive-definite input
$L(P_{e_0})=36I_4$.  Proposition~\ref{prop:twenty-directions} supplies
the twenty distinct directions and the quadratic bound.  Proposition
\ref{prop:tensor-criterion} therefore applies, with the explicit maps
in \eqref{eq:phi-pair} determined by the integer matrices
\eqref{eq:explicit-matrices}.

The positive test also gives an exact nonzero Choi value.  Since
$\widehat{e_0}=e_{00}$, Equations~\eqref{eq:SR} and
\eqref{eq:choi-pairing} yield
\[
 S(P_{e_{00}})=R(P_{e_{00}})=36^2I_6,\qquad
 (e_{00}\otimes e_{00})^*Z(e_{00}\otimes e_{00})
 =6\cdot36^4=10{,}077{,}696>0.
\]
Thus the constructed pair has the asserted nonzero product-free
composite Choi range, and its composite is not entanglement breaking.
\end{proof}
